\section*{Abstract}

Network-proximity Z-scores are often compared across candidate compounds, but unequal target counts change the null variance of the averaged distance statistic. In a liver-expressed STRING interactome, we show that the null standard deviation shrinks approximately as $|T|^{-1/2}$, so standardized evidence can rank a larger target set above a smaller set even when its raw closest-distance effect is weaker. Hyperforin (10 targets) is closer to the drug-induced liver injury (DILI) module than Quercetin (62 targets) ($d_c = 1.30$ vs $1.68$, STRING $\geq$900) yet has the weaker proximity Z-score ($-3.86$ vs $-5.44$); the Z-score ordering also reverses between STRING thresholds while the raw-distance ordering does not. A degree-controlled calibration benchmark reproduces the same null-precision law across the DILI module and matched pseudo-modules, placing material rank reversal in a narrow, high-margin regime rather than as a generic outcome. As a complementary effect-size readout, perturbation efficiency (mean per-target influence) separates direct target--DILI overlap from propagated random-walk influence: 62\% of Hyperforin's efficiency is direct overlap, leaving a modest propagated advantage ($\sim$1.5$\times$). Target-asymmetric network comparisons should therefore report raw effect size, standardized evidence, and direct versus propagated influence separately.

\vspace{0.5em}
\noindent\textbf{Keywords:} network medicine; network proximity; proximity Z-score; target-count asymmetry; effect size; random walk with restart; drug-induced liver injury.
